\chapter{Design Decisions}\label{r02:ch:design}

The gaps of \cref{r02:ch:gap} could be closed in many ways. This chapter explains the
choices that were made, and why a first, more ambitious implementation was replaced by
a much smaller one.

\section{Two attempts}\label{r02:sec:design-history}

\paragraph{First attempt: VLS and VLA, configurable (October 2026).}
The first implementation (\ps{} commit \texttt{31758e93}, waLBerla \texttt{e9712d59})
aimed at full generality. In \ps{} it added about 2600 lines:
\begin{itemize}
  \item both vector-length-specific and vector-length-agnostic (stripmining) code, with
        a new IR function \code{PsVecStripmineLanes} and a new \code{stripmine} loop
        transformation;
  \item a configurable VLEN, a register-group multiplier (LMUL) derived per type from
        lanes and VLEN, and new configuration classes (\code{RvvOptions},
        \code{VectorizationMode});
  \item an optional run-time VLEN check in every kernel, broadcast folding, \code{vid}
        and slide chains for non-uniform constants, and \code{-march} synthesis from
        \file{/proc/cpuinfo}.
\end{itemize}
In waLBerla it added three CMake options and plumbed them into sweepgen
(\cref{r02:lst:first-attempt}).

\rneedspace{12\baselineskip}
\begin{diffcode}[nolol=false,caption={[The first attempt in waLBerla]The first attempt in waLBerla (commit \texttt{e9712d59}, abridged): new options, mapped into sweepgen, which then forced vectorization.},label=r02:lst:first-attempt]
--- CMakeLists.txt
+ option ( WALBERLA_RISCV_RVV  "Generate kernels for the RISC-V V extension (RVV 1.0)" OFF )
+ set    ( WALBERLA_RVV_VLEN "128" CACHE STRING "Assumed (VLS) / minimum (VLA) RVV VLEN in bits" )
+ set    ( WALBERLA_RVV_MODE "vla" CACHE STRING "RVV vectorization mode: 'fixed' (VLS) or 'vla'" )
--- sweepgen/src/sweepgen/build_config.py
+        elif self.rvv_enabled:
+            target = Target.RISCV_RVV
 ...
+        if cfg.target == Target.RISCV_RVV:
+            cfg.cpu.vectorize.enable = True
+            cfg.cpu.vectorize.assume_inner_stride_one = True
+            cfg.cpu.vectorize.mode = self.rvv_mode
+            cfg.cpu.rvv.vlen = self.rvv_vlen
\end{diffcode}

It worked: about 150 tests passed on the board, and kernels gave bit-identical results
for VLEN 128 to 1024 under QEMU. But it was large, it introduced concepts no other
\ps{} backend has, and it made RVV builds of waLBerla behave differently from every
other CPU build (vectorization forced on, new options to learn).

\paragraph{Second attempt: VLS-256 with AVX parity (November 2026).}
The current implementation (\ps{} \texttt{9a989e16}, waLBerla \texttt{dbb1b400})
started by restoring every file of the first attempt to upstream, and then rebuilt the
backend under one rule: \emph{make the same assumptions the AVX backend makes}. The
\ps{} diff shrank from about 2600 to about 760 lines of source (1270 with tests); the
waLBerla build change shrank to about 20 lines of CMake. On the LBM workloads of the
first attempt, the simpler backend measured the same within $\pm 2$\,\% in double
precision and within run-to-run noise in single precision. The plan of this rewrite is
\file{docs/notes2/rvv-vls256-only-plan-opus5.5.md}.

\section{Vector-length specific, VLEN = 256}\label{r02:sec:design-vls}

Fixing VLEN turns RVV into a fixed-width SIMD instruction set, so the whole machinery
\ps{} already has for AVX applies unchanged: lane counts known at code generation,
peeling, a vector loop and a scalar remainder, one register type per vector width.
\begin{description}[style=nextline]
  \item[What it costs] Generated code is correct only on hardware with VLEN $\geq$ 256,
    and uses only 256 bits of wider registers. The K1 is the only RVV hardware in this
    project, so neither costs anything today.
  \item[What it saves] No new IR node, no new loop strategy, no new configuration
    options. The backend is a platform class plus a runtime header, like the x86 one.
  \item[LMUL] Only one register (\code{m1}) and half a register (\code{mf2}, needed for
    mixed-precision casts, mirroring AVX's \code{__m128}). Register groups
    (\code{m2}, \code{m4}) behave like wider-than-native vectors; AVX rejects those too.
    An \code{m2} variant measured faster for one f64 collide kernel during the first
    attempt; adding it later is a matter of the width table in \code{_lmul}.
  \item[vl] Every intrinsic receives the literal lane count (\code{4uLL} for
    \code{double}). At VLEN = 256 that equals VLMAX for every legal type, so each
    operation processes exactly one full (or half) register, and no \code{vsetvl}
    bookkeeping is visible in the generated code. The compiler inserts the
    \code{vsetvli} instructions.
\end{description}

\section{No new waLBerla options}\label{r02:sec:design-no-options}

On an x86 workstation, a developer who wants host-specific code sets
\cmakeopt{WALBERLA_OPTIMIZE_FOR_LOCALHOST=ON}. The compiler gets \code{-march=native},
sweepgen selects \code{CurrentCPU}, and \ps{} resolves it to the best x86 target. The
second attempt keeps exactly this path for RISC-V. Its two halves only had to be made to
work:
\begin{itemize}
  \item the compiler half, by replacing \code{-march=native} with an explicit string on
        riscv64 (\cref{r02:sec:wlb-cmake});
  \item the code generator half, by teaching \ps{} to detect RVV and resolve
        \code{CurrentCPU} to \code{RISCV_RVV} (\cref{r02:sec:ps-target}).
\end{itemize}
The first attempt's options \cmakeopt{WALBERLA_RISCV_RVV}, \code{_VLEN} and \code{_MODE}
became unnecessary once VLEN was fixed and VLA dropped. Equally, sweepgen no longer
forces vectorization for RVV. Whether a sweep is vectorized is the generator script's
decision on every platform; making RVV the only platform where it happens automatically
would surprise anyone who knows waLBerla on x86.

\section{Compiler flags: \texttt{zvl256b}, not \texttt{-mrvv-vector-bits}}\label{r02:sec:design-flags}

Both waLBerla and the \ps{} JIT use \code{-march=rv64gcv_zvl256b}:
\begin{description}
  \item[\code{rv64gc}] the base ISA with the usual extensions (what the compilers use
    by default on the board);
  \item[\code{v}] the vector extension, required for any RVV intrinsic;
  \item[\code{zvl256b}] the promise ``VLEN $\geq$ 256''. It makes the compiler define
    \code{__riscv_v_min_vlen = 256}, which the runtime header checks
    (\cref{r02:sec:design-guard}).
\end{description}
GCC and Clang also accept \code{-mrvv-vector-bits=zvl}, which makes the compiler assume
VLEN \emph{equal} to the \code{Zvl} value and treat the vector types as fixed-size. The
backend does not need that: its \code{vl} arguments are literals anyway. It would also
narrow the contract from ``VLEN $\geq$ 256'' to ``VLEN = 256''. The first attempt
measured it as performance-neutral, so it was left out.

The same string is produced in two places that must agree: waLBerla's CMake files for
the C++ build, and \ps{}' \code{CompilerInfo} for the JIT. Both apply the same rule to
the same input: if the \code{isa} line of \file{/proc/cpuinfo} lists the single-letter
extension \code{v}, use \code{-march=rv64gcv_zvl256b}, otherwise \code{-march=rv64gc}.

\section{A compile-time guard instead of a run-time check}\label{r02:sec:design-guard}

The AVX backend has no run-time CPU check: a kernel compiled with
\code{-march=x86-64-v3} simply crashes with an illegal instruction on a CPU without
AVX2. The RVV backend follows suit, with one addition: the failure mode on RVV hardware
with VLEN $<$ 256 is \emph{silent} (fewer lanes processed, wrong results), not a crash.
The runtime header therefore refuses to compile unless the compiler was told
VLEN $\geq$ 256:
\begin{ccode}
#if !defined(__riscv_v_min_vlen) || __riscv_v_min_vlen < 256
#error "pystencils RVV kernels assume VLEN >= 256 bits: compile with -march=rv64gcv_zvl256b"
#endif
\end{ccode}
This catches wrong build flags. It cannot catch a binary built correctly and then run on
a smaller-VLEN machine; that remains a documented precondition (\cref{r02:sec:lim-vlen}).
The first attempt emitted a run-time guard at the entry of every vector kernel: it read
the hardware VLEN with \code{__riscv_vlenb()} and called \code{std::abort()} if it did
not match. It was dropped for parity with AVX.

\section{Detecting RVV}\label{r02:sec:design-detect}

Linux reports the ISA string in \file{/proc/cpuinfo}, e.g.\ on the K1:
\begin{plaincode}
isa : rv64imafdcv_zicbom_zicboz_..._zve32f_zve32x_zve64d_zve64f_zve64x_zvfh_zvfhmin_...
\end{plaincode}
RVV~1.0 appears as the single letter \code{v} in the leading part before the first
underscore. Detection looks only there, on purpose: multi-letter extensions such as
\code{zve64d} (embedded vector subsets) or \code{xtheadvector} (the incompatible RVV~0.7
draft of T-Head cores) must not count. The CMake files use the regular expression
\code{rv64[a-z]*v} on the same line, which also stops at the first underscore.

\section{Summary}

\begin{center}\small
\begin{tabularx}{\linewidth}{@{}>{\raggedright\arraybackslash}p{4.6cm}>{\raggedright\arraybackslash}X@{}}
  \toprule
  \textbf{Decision} & \textbf{Reason}\\
  \midrule
  VLS with VLEN = 256, LMUL $\leq$ 1 & reuses the AVX code path; the K1 is the only
    target; measured no slower than the general version\\
  no new waLBerla CMake options & \cmakeopt{WALBERLA_OPTIMIZE_FOR_LOCALHOST} already
    means ``this host''; same behaviour as x86\\
  no sweepgen change & target selection is generic; vectorization is opt-in on every
    platform\\
  \code{-march=rv64gcv_zvl256b} & enables V and guarantees VLEN $\geq$ 256 to the
    compiler and the runtime header\\
  no \code{-mrvv-vector-bits} & not needed by literal \code{vl}; would narrow the
    contract; performance-neutral\\
  compile-time \texttt{\#error}, no run-time check & catches wrong flags; parity with
    AVX, which has no run-time check either\\
  detect single-letter \code{v} only & exclude \code{zve*} and RVV~0.7
    (\code{xtheadvector})\\
  \bottomrule
\end{tabularx}
\end{center}
